\documentclass{report}
\usepackage{luatexja} % LuaTeXで日本語を使うためのパッケージ
\usepackage{luatexja-fontspec} % LuaTeX用の日本語フォント設定

% --- 数学関連 ---
\usepackage{amsmath, amssymb, amsfonts, mathtools, bm, amsthm} % 基本的な数学パッケージ
\usepackage{type1cm, upgreek} % 数式フォントとギリシャ文字
\usepackage{physics, mhchem} % 物理や化学の記号や式の表記を簡単にする

% --- 表関連 ---
\usepackage{multirow, longtable, tabularx, array, colortbl, dcolumn, diagbox} % 表のレイアウトを柔軟にする
\usepackage{tablefootnote, truthtable} % 表中に注釈を追加、真理値表
\usepackage{tabularray} % 高度な表組みレイアウト

% --- グラフィック関連 ---
\usepackage{tikz, graphicx} % 図の描画と画像の挿入
\usepackage{background} % ウォーターマークの設定
\usepackage{caption, subcaption} % 図や表のキャプション設定
\usepackage{float, here} % 図や表の位置指定

% --- レイアウトとページ設定 ---
\usepackage{fancyhdr} % ページヘッダー、フッター、余白の設定
\usepackage[top = 20truemm, bottom = 20truemm, left = 20truemm, right = 20truemm]{geometry}
\usepackage{fancybox, ascmac} % ボックスのデザイン

% --- 色とスタイル ---
\usepackage{xcolor, color, colortbl, tcolorbox} % 色とカラーボックス
\usepackage{listings, jvlisting} % コードの色付けとフォーマット

% --- 参考文献関連 ---
\usepackage{biblatex, usebib} % 参考文献の管理と挿入
\usepackage{url, hyperref} % URLとリンクの設定

% --- その他の便利なパッケージ ---
\usepackage{footmisc} % 脚注のカスタマイズ
\usepackage{multicol} % 複数段組
\usepackage{comment} % コメントアウトの拡張
\usepackage{siunitx} % 単位の表記
\usepackage{docmute}
% \usepackage{appendix}
% --- tcolorboxとtikzの設定 ---
\tcbuselibrary{theorems, breakable} % 定理のボックスと改ページ設定
\usetikzlibrary{decorations.markings, arrows.meta, calc} % tikzの装飾や矢印の設定

% --- 定理スタイルと数式設定 ---
\theoremstyle{definition} % 定義スタイル
\numberwithin{equation}{section} % 式番号をサブセクション単位でリセット

% --- hyperrefの設定 ---
\hypersetup{
  setpagesize = false,
  bookmarks = true,
  bookmarksdepth = tocdepth,
  bookmarksnumbered = true,
  colorlinks = false,
  pdftitle = {}, % PDFタイトル
  pdfsubject = {}, % PDFサブジェクト
  pdfauthor = {}, % PDF作者
  pdfkeywords = {} % PDFキーワード
}

% --- siunitxの設定 ---
\sisetup{
  table-format = 1.5, % 小数点以下の桁数
  table-number-alignment = center, % 数値の中央揃え
}

% --- 透かし画像の設定 --- 
\backgroundsetup{
  scale=0.5,                       % 画像のスケール
  % color=black,                   % 画像の色（透かし用に半透明が推奨）
  opacity=0.2,                   % 透かしの透明度（0が完全透明、1が完全不透明）
  angle=0,                       % 画像の角度
  position = current page.south east,  % ページの右下
  hshift=-6cm, % 右方向へのシフト（負の値で内側に移動）
  vshift=5cm,  % 上方向へのシフト（正の値で内側に移動）
  contents={\includegraphics{./fig/appilogo-circular-full.png}} % 画像のパス
}

% --- その他の設定 ---
\allowdisplaybreaks % 数式の途中改ページ許可
\newcolumntype{t}{!{\vrule width 0.1pt}} % 新しいカラムタイプ
\newcolumntype{b}{!{\vrule width 1.5pt}} % 太いカラム
\UseTblrLibrary{amsmath, booktabs, counter, diagbox, functional, hook, html, nameref, siunitx, varwidth, zref} % tabularrayのライブラリ
\setlength{\columnseprule}{0.4pt} % カラム区切り線の太さ
\captionsetup[figure]{font = bf} % 図のキャプションの太字設定
\captionsetup[table]{font = bf} % 表のキャプションの太字設定
\captionsetup[lstlisting]{font = bf} % コードのキャプションの太字設定
\captionsetup[subfigure]{font = bf, labelformat = simple} % サブ図のキャプション設定
\setcounter{secnumdepth}{4} % セクションの深さ設定
\newcolumntype{d}{D{.}{.}{5}} % 数値のカラム
\newcolumntype{M}[1]{>{\centering\arraybackslash}m{#1}} % センター揃えのカラム
\DeclareMathOperator{\diag}{diag}
\everymath{\displaystyle} % 数式のスタイル
\newcommand{\inner}[2]{\left\langle #1, #2 \right\rangle}
\renewcommand{\figurename}{図}
\renewcommand{\i}{\mathrm{i}} % 複素数単位i
\renewcommand{\laplacian}{\grad^2} % ラプラシアンの記号
\renewcommand{\thesubfigure}{(\alph{subfigure})} % サブ図の番号形式
\newcommand{\m}[3]{\multicolumn{#1}{#2}{#3}} % マルチカラムのショートカット
\renewcommand{\r}[1]{\mathrm{#1}} % mathrmのショートカット
\newcommand{\e}{\mathrm{e}} % 自然対数の底e
\newcommand{\Ef}{E_{\mathrm{F}}} % フェルミエネルギー
\renewcommand{\c}{\si{\degreeCelsius}} % 摂氏記号
\renewcommand{\d}{\r{d}} % d記号
\renewcommand{\t}[1]{\texttt{#1}} % タイプライタフォント
\newcommand{\kb}{k_{\mathrm{B}}} % ボルツマン定数
% \renewcommand{\phi}{\varphi} % ϕをφに変更
\renewcommand{\epsilon}{\varepsilon}
\newcommand{\fullref}[1]{\textbf{\ref{#1} \nameref{#1}}}
\newcommand{\reff}[1]{\textbf{図\ref{#1}}} % 図参照のショートカット
\newcommand{\reft}[1]{\textbf{表\ref{#1}}} % 表参照のショートカット
\newcommand{\refe}[1]{\textbf{式\eqref{#1}}} % 式参照のショートカット
\newcommand{\refp}[1]{\textbf{コード\ref{#1}}} % コード参照のショートカット
\renewcommand{\lstlistingname}{コード} % コードリストの名前
\renewcommand{\theequation}{\thesection.\arabic{equation}} % 式番号の形式
\renewcommand{\footrulewidth}{0.4pt} % フッターの線
\newcommand{\mar}[1]{\textcircled{\scriptsize #1}} % 丸囲み文字
\newcommand{\combination}[2]{{}_{#1} \mathrm{C}_{#2}} % 組み合わせ
\newcommand{\thline}{\noalign{\hrule height 0.1pt}} % 細い横線
\newcommand{\bhline}{\noalign{\hrule height 1.5pt}} % 太い横線

% --- カスタム色定義 ---
\definecolor{burgundy}{rgb}{0.5, 0.0, 0.13} % バーガンディ色
\definecolor{charcoal}{rgb}{0.21, 0.27, 0.31} % チャコール色
\definecolor{forest}{rgb}{0.0, 0.35, 0} % 森の緑色

% --- カスタム定理環境の定義 ---
\newtcbtheorem[number within = chapter]{myexc}{練習問題}{
  fonttitle = \gtfamily\sffamily\bfseries\upshape,
  colframe = forest,
  colback = forest!2!white,
  rightrule = 1pt,
  leftrule = 1pt,
  bottomrule = 2pt,
  colbacktitle = forest,
  theorem style = standard,
  breakable,
  arc = 0pt,
}{exc-ref}
\newtcbtheorem[number within = chapter]{myprop}{命題}{
  fonttitle = \gtfamily\sffamily\bfseries\upshape,
  colframe = blue!50!black,
  colback = blue!50!black!2!white,
  rightrule = 1pt,
  leftrule = 1pt,
  bottomrule = 2pt,
  colbacktitle = blue!50!black,
  theorem style = standard,
  breakable,
  arc = 0pt
}{proposition-ref}
\newtcbtheorem[number within = chapter]{myrem}{注意}{
  fonttitle = \gtfamily\sffamily\bfseries\upshape,
  colframe = yellow!20!black,
  colback = yellow!50,
  rightrule = 1pt,
  leftrule = 1pt,
  bottomrule = 2pt,
  colbacktitle = yellow!20!black,
  theorem style = standard,
  breakable,
  arc = 0pt
}{remark-ref}
\newtcbtheorem[number within = chapter]{myex}{例題}{
  fonttitle = \gtfamily\sffamily\bfseries\upshape,
  colframe = black,
  colback = white,
  rightrule = 1pt,
  leftrule = 1pt,
  bottomrule = 2pt,
  colbacktitle = black,
  theorem style = standard,
  breakable,
  arc = 0pt
}{example-ref}
\newtcbtheorem[number within = chapter]{exc}{Requirement}{myexc}{exc-ref}
\newcommand{\rqref}[1]{{\bfseries\sffamily 練習問題 \ref{exc-ref:#1}}}
\newtcbtheorem[number within = chapter]{definition}{Definition}{mydef}{definition-ref}
\newcommand{\dfref}[1]{{\bfseries\sffamily 定義 \ref{definition-ref:#1}}}
\newtcbtheorem[number within = chapter]{prop}{命題}{myprop}{proposition-ref}
\newcommand{\prref}[1]{{\bfseries\sffamily 命題 \ref{proposition-ref:#1}}}
\newtcbtheorem[number within = chapter]{rem}{注意}{myrem}{remark-ref}
\newcommand{\rmref}[1]{{\bfseries\sffamily 注意 \ref{remark-ref:#1}}}
\newtcbtheorem[number within = chapter]{ex}{例題}{myex}{example-ref}
\newcommand{\exref}[1]{{\bfseries\sffamily 例題 \ref{example-ref:#1}}}
% --- 再定義コマンド ---
% \mathtoolsset{showonlyrefs=true} % 必要な式番号のみ表示
\pagestyle{fancy} % ヘッダー・フッターのスタイル設定
\chead{応用量子物性講義ノート} % 中央ヘッダー
% \rhead{}
\fancyhead[R]{\rightmark}
\renewcommand{\sectionmark}[1]{\markright{\thesection\ #1}}
\cfoot{\thepage} % 中央フッターにページ番号
\lhead{}
\rfoot{Yuto Masuda and Haruki Aoki} % 右フッターに名前
\setcounter{tocdepth}{4} % 目次の深さ
\makeatletter
\@addtoreset{equation}{section} % サブセクションごとに式番号をリセット
\makeatother

% --- メタ情報 ---
\title{応用量子物性講義ノート}
\date{更新日\today}
\author{Yuto Masuda and Haruki Aoki}

\begin{document}
  この章では解析力学から量子力学への接続を駆け足で書く．
  座標と運動量に依存するある物理量$\omega(x_i,p_i)$を考える．$\omega$の時間微分は
  \begin{align}
    \dv{\omega}{t} = \frac{\partial\omega}{\partial x_i}\dot{x}_i + \frac{\partial \omega}{\partial p_i}\dot{p}_i \label{time-derivative-omega}
  \end{align}
  である．ここで正準方程式
  \begin{align}
    \begin{dcases}
      \dot{x}_i = \frac{\partial H}{\partial p_i}\\
      \dot{p}_i = -\frac{\partial H}{\partial x_i}
    \end{dcases}
  \end{align}
  より，\refe{time-derivative-omega}の右辺は
  \begin{align}
    \frac{\partial\omega}{\partial x_i}\dot{x}_i + \frac{\partial \omega}{\partial p_i}\dot{p}_i
    =
    \frac{\partial\omega}{\partial x_i}\frac{\partial H}{\partial p_i} - \frac{\partial \omega}{\partial p_i}\frac{\partial H}{\partial x_i}
  \end{align}
  と書ける．上式の右辺を
  \begin{align}
    \qty{\omega,H} \equiv   \frac{\partial\omega}{\partial x_i}\frac{\partial H}{\partial p_i} - \frac{\partial \omega}{\partial p_i}\frac{\partial H}{\partial x_i}
  \end{align}
  とおき，これを\textbf{Poisson括弧}という．つまり，$\omega$の時間微分は
  \begin{align}
    \dv{\omega}{t} = \qty{\omega,H} \label{poisson-h}
  \end{align}
  と書ける．Poisson括弧は一般の物理量$\lambda$で
  \begin{align}
    \qty{\omega,\lambda} \equiv   \frac{\partial\omega}{\partial x_i}\frac{\partial \lambda}{\partial p_i} - \frac{\partial \omega}{\partial p_i}\frac{\partial \lambda}{\partial x_i}
  \end{align}
  と定義される．

  $\omega$の時間発展
  \begin{align}
    \omega(t) = \omega(t+\epsilon) = \omega(t) + \epsilon\dv{\omega}{t}
  \end{align}
  を考える．これは\refe{poisson-h}を用いると，
  \begin{align}
    \omega(t+\epsilon) = \omega(t) + \epsilon\qty{\omega,H}
  \end{align}
  と書くことができる．つまり，$\omega$の時間発展による変化量$\delta\omega=\epsilon\dv{\omega}{t}$はPoisson括弧により，
  \begin{align}
    \delta\omega = \epsilon\qty{\omega,H} \label{time-evolution-by-poisson}
  \end{align}
  と表される．\refe{time-evolution-by-poisson}から，\textbf{$\omega$の時間発展が$\omega$と$H$のPoisson括弧をとることにより生み出される}ことがわかる．
  もし$\omega$が保存量であれば
  \begin{align}
    \qty{\omega,H}=0
  \end{align}
  である．

  この概念を一般化する．物理量$\omega$が何らかの変換により，
  \begin{align}
    \omega \to \omega + \delta\omega = \omega + \epsilon\qty[    \frac{\partial\omega}{\partial x_i}\delta x_i + \frac{\partial \omega}{\partial p_i}\delta p_i
    ]
  \end{align}
  と変化したとする．ここである関数$Q(x_i,p_i)$が$\delta x_i,\delta p_i$を生み出したとする．つまり，
  \begin{align}
    \delta x_i &= \frac{\partial Q}{\partial p_i}\\
    \delta p_i &= -\frac{\partial Q}{\partial x_i}
  \end{align}
  とすると，
  $\omega$の変化は，
  \begin{align}
    \omega + \epsilon\qty{\omega,Q}
  \end{align}
  と書ける．$Q$はPoisson括弧をとることにより変換を生み出したとみなせるので，$Q$は\textbf{Generator}と呼ばれる．
  \begin{myex}{}{}
  $\omega=x_i$，$Q=p_j$とする．
  \begin{align}
    \qty{x_i, p_j} &= \frac{\partial\x_i}{\partial x_i}\frac{\partial p_j}{\partial p_i} + \frac{\partial \x_i}{\partial p_i}\frac{\partial p_i}{\partial x_i}\\
    &= \delta_{ji}
  \end{align}
  よって，
  \begin{align}
    x_i \to x_i + \delta x_i = x_i + \epsilon \delta_{ji}
  \end{align}
  したがって，$j$方向の運動量は$j$方向への並進操作を生成する．
  \end{myex}
  まとめると，$H$のPoisson括弧は時間発展，$p$のPoisson括弧は並進操作を生み出す．同様に角運動量$l$のPoisson括弧が回転操作に対応することも
  確認できる．また，$x$と$p$の間には
  \begin{align}
    \qty{x,p} = 1 \label{xp-poisson}
  \end{align}
  が成り立つ．

  量子力学の世界に入る．$p$は並進操作を生むので，運動量演算子$\hat{p}$が並進操作のgeneratorと考える．
  並進操作の演算子を
  \begin{align}
    \e^{-\i\frac{a\hat{p}}{\hbar}}\ket{x} = \ket{x+a}
  \end{align}
  と定義する．$a$が微小量$\epsilon$のとき，
  \begin{align}
    \qty(\hat{I} - \frac{\i}{\hbar}\epsilon\hat{p}) \ket{x} = \ket{x+\epsilon}
  \end{align}
  と書ける．これを用いて座標空間での$\hat{p}$の形を導出する．ここで，座標空間では
  \begin{align}
    \hat{x}\ket{\psi} \to x\psi(x)
  \end{align}
  である．

  一般の量子状態$\ket{\psi}$に$\ket{x+\epsilon}$を作用させて，
  \begin{align}
    \bra{x+\epsilon}\ket{\psi} &= \bra{x}\qty(\hat{I} + \frac{\i}{\hbar}\epsilon\hat{p})\ket{\psi}\\
    &=\braket{x}{\psi} + \frac{\i}{\hbar}\epsilon\bra{x}\hat{p}\ket{\psi}\\
    &= \braket{x}{\psi} + \frac{\i}{\hbar} \epsilon\int\dd{x'}\bra{x}\hat{p}\ket{x'}\bra{x'}\ket{\psi}
  \end{align}
  を得る．
  また，$\braket{x+\epsilon}{\psi}$は波動関数$\psi(x)$を用いて，
  \begin{align}
    \braket{x+\epsilon}{\psi} &= \psi(x+\epsilon)\\
    &= \psi(x) + \epsilon\dv{\psi}{x}\\
    &= \psi(x) + \epsilon \int \dd{x'}\delta(x-x')\dv{\psi(x')}{x'}\\
    &= \psi(x) - \epsilon\int\dd{x'}\dv{\delta(x-x')}{x'}\psi(x')
  \end{align}
  とも書ける．これら2式を見比べると，
  \begin{align}
    \bra{x}\hat{p}\ket{x'} = -\i\hbar\dv{\delta(x-x')}{x}
  \end{align}
  が得られる．よって，
  \begin{align}
    \bra{x}\hat{p}\ket{\psi} = -\i\hbar\dv{\psi}{x}
  \end{align}
  である．したがって，
  \begin{align}
    \hat{p} \to -\i\hbar\frac{\partial}{\partial x}
  \end{align}
  とできる．さらに，$\hat{x}$と$\hat{p}$の間には交換関係
  \begin{align}
    [\hat{x},\hat{p}] = \i\hbar
  \end{align}
  が成り立つ．これは解析力学でのPoisson括弧\refe{xp-poisson}の$\i\hbar$倍である．
\end{document}
